\subsection{INFINITE SUMS OF COMPLEX NUMBERS}

Since the additive group of the field C is the same as the additive group \(R^{2}\) , it is not necessary to make a special study of summable families and series in C, since this is included in the general theory of Chapter VII, § 3; we leave to the reader the exercise of translating the results of this theory into the language of the theory of complex numbers. We state only the following proposition, which is a corollary to Proposition 3 of Chapter VII, § 3, no. 1:

PROPOSITION 1. If \((u_{\lambda})_{\lambda \in \mathbf{L}}\) and \((v_{\mu})_{\mu \in \mathbf{M}}\) are two summable families of complex numbers, then the family \((u_{\lambda}v_{\mu})_{(\lambda ,\mu)\in \mathbf{L}\times \mathbf{M}}\) is summable and we have

\[
\sum_ {(\lambda , \mu) \in \mathbf {L} \times \mathbf {M}} u _ {\lambda} v _ {\mu} = \left(\sum_ {\lambda \in \mathbf {L}} u _ {\lambda}\right) \left(\sum_ {\mu \in \mathbf {M}} v _ {\mu}\right).\tag{1}
\]

We leave it to the reader to state the corresponding result for quaternions.

